\chapter{基于模型的知识获取}
\label{chap:knowledge-acquisition-model-based}

一批花卉照片已经采集完成，却还没有逐株对应的物种答案。人工可以逐一判断，规则可以在叶缘、花期等条件明确时给出候选；照片很多、可见部位各不相同，或鉴别依据难以穷尽成规则时，我们还可以让模型参与作答。模型可以利用过去数据中学到的关系，也可以在调用时读取带答案的参考照片或样本间的连接，形成类别、数值、对象关系或文本答案。这些输出形成记录后，沿用第~\ref{chap:knowledge-acquisition-overview}章的名称，称为\term{候选知识}：它们已经取得，但能否作为当前任务的依据或训练材料，仍待检查。

调用完成只回答了“模型返回了什么”，还没有回答“这条输出依据什么、适用于哪里”。要继续使用答案，先要确认模型能接收当前照片并按所用物种表输出，再确认调用中实际提供了哪些观测、示范和外部材料，随后辨认答案是由参数映射、参考样本还是图上连接算出的。错误反馈也须沿这些依赖定位，才能决定改模型、改输入还是改下一轮的处理对象。本章要解决的就是：怎样把准备、调用、预测和反馈连起来，获得可追溯的任务候选，同时保留不参与方案选择的独立核验依据？

在贯穿本章的花卉任务中，我们要为照片补充物种名称；证据不足时允许弃权，并保存判断依据。可用资源可能是少量已标注照片、一个现成识别模型，或能接收图文指令的通用模型。资源不同，准备工作会改变；记录要求不变：每条输出都要对应具体模型版本、输入版本和预测机制，复核者才能追查错误在哪一步进入。

下面沿这条依赖关系推进：先得到可调用的模型，再为它组织当前输入，由预测机制产生候选记录，最后根据复核结果选择下一轮的修改。记录仍按定义~\ref{def:knowledge-acquisition-record}分开保存内容、来源动作、执行状态与验证状态；一次成功的弃权不等于调用失败，一条已经给出的物种名也不等于已经核验。

\section{标注模型的准备}
\label{sec:knowledge-acquisition-model-preparation}

面对尚无物种答案的照片，第一步是找出能按本任务要求作答的模型。本章把承担这一工作的模型称为\term{标注模型}（labeling model），也就是为待处理对象提供候选知识的模型，不特指某种网络。准备时，我们先从可用资源中选出可试验方案；需要改参数的方案再经过训练或适配；得到可调用版本后，用少量照片试运行，确定已经检查的能力范围。这三步分别交付候选方案、模型版本和采用或继续修改的依据。

是否训练，要看现有模型与任务之间缺了什么，不能只看模型来源。专用识别模型若使用另一套物种表，可能需要适配；通用模型若已能按任务说明和示范作答，可能无需更新参数；从头构建则可能因标签不足而暂不能训练。沿图~\ref{fig:knowledge-acquisition-model-preparation}的分支可以看到，不需要训练的方案直接进入试标。试标后若问题在接口、标签或任务范围，反馈就返回相应准备步骤。

\begin{figure}[htbp]
  \centering
  % 图4-1：标注模型的准备与试标决策，按112 mm正文宽度设计。
\begin{tikzpicture}[
  x=1mm,
  y=1mm,
  font=\figurefont\footnotesize,
  text=BookInk,
  prep input/.style={book node,draw=FigureInput,fill=FigureInput!7,line width=0.8pt,text width=25mm,minimum height=11mm,inner sep=1.5mm},
  prep process/.style={book node,draw=FigureModel,fill=FigureModel!7,line width=0.8pt,text width=28mm,minimum height=11mm,inner sep=1.5mm},
  prep decision/.style={book node,draw=BookAmber,fill=BookAmber!7,line width=0.9pt,text width=24mm,minimum height=11mm,inner sep=1.5mm},
  prep result/.style={book node,draw=BookMuted,fill=BookPaper,line width=0.7pt,text width=25mm,minimum height=11mm,inner sep=1.5mm},
  prep flow/.style={knowledge arrow,draw=BookMuted,line width=0.8pt,rounded corners=1mm},
  prep feedback/.style={knowledge arrow,draw=BookAmber,line width=0.8pt,dashed,rounded corners=1mm},
  edge text/.style={fill=white,inner sep=0.8pt,font=\figurefont\scriptsize,text=BookMuted}
]
  \node[prep input] (requirements) at (16,0) {任务要求\\类别、边界、弃权};
  \node[prep input] (models) at (55,0) {可用模型\\专用或通用};
  \node[prep input] (labels) at (94,0) {初始标签\\样本与覆盖};

  \node[prep process] (select) at (55,-19) {选择候选方案\\并记录资源约束};
  \node[prep decision] (need) at (55,-37) {是否需要\\模型参数训练？};
  \node[prep process] (adapt) at (87,-51) {训练或适配\\形成模型版本};
  \node[prep process] (pilot) at (55,-66) {试标确认\\能力范围与失败类型};

  \node[prep result,draw=BookAmber] (source) at (17,-86) {补充其他来源\\或重定任务};
  \node[prep result,draw=FigureOutput] (accept) at (55,-86) {采用当前方案\\进入批量获取};
  \node[prep result,draw=FigureModel] (continue) at (93,-86) {继续适配\\再做试标};

  \draw[prep flow] (requirements.south) -- (16,-10) -| (select.west);
  \draw[prep flow] (models.south) -- (select.north);
  \draw[prep flow] (labels.south) -- (94,-10) -| (select.east);
  \draw[prep flow] (select.south) -- (need.north);
  \draw[prep flow] (need.south) -- node[edge text,right] {否} (pilot.north);
  \draw[prep flow] (need.east) -- node[edge text,above] {是} (72,-37) |- (adapt.north);
  \draw[prep flow] (adapt.south) |- (pilot.east);
  \draw[prep flow] (pilot.south) -- (accept.north);
  \draw[prep flow] (pilot.west) -- (36,-66) |- (source.north);
  \draw[prep flow] (pilot.east) -- (75,-66) |- (continue.north);

  \draw[prep feedback] (source.west) -- (1,-86) |- (select.west);
  \draw[prep feedback] (continue.east) -- (110,-86) |- (adapt.east);
  \node[edge text,anchor=west] at (3,-29) {资源或要求变化};
  \node[edge text,anchor=east] at (108,-42) {新训练版本};
\end{tikzpicture}

  \caption{标注模型的准备与试标决策。上方三类资源共同限定候选方案；“是否需要参数训练”分支决定直接试标还是先适配。下方反馈返回需要改变的步骤。试标用于准备阶段的方案选择，不能代替方案冻结后的独立验收。}
  \label{fig:knowledge-acquisition-model-preparation}
\end{figure}

\subsection{模型方案的选择}
\label{sec:knowledge-acquisition-model-scheme-selection}

要比较可用模型，先把花卉任务的输入、输出和范围定下来。输入可仅含照片，也可包括地点、季节和植株部位；输出可要求物种表中的一个名称或多个候选，同时约定证据不足时如何返回未知并弃权。覆盖本地常见花卉与不预先限定物种的植物照片，需要检查的类别范围也不同。这些条件确定后，才有依据判断现有模型是否支持所需信息类型、类别含义和输出字段。

在这些共同要求下，可以按需要完成的准备动作比较三种方案。没有可用已有模型而有任务数据和标签时，从头构建是一个候选；已有模型的任务语义与输入条件相容时，可以试验直接复用；预训练参数有可用表示，但类别、数据分布或接口不完全一致时，可以试验适配。适配是让已有模型按当前任务工作，不保证少量更新就能弥合任何差异。专用模型和通用模型都可能直接复用或需要适配，模型大小也不说明它已经适用。模型卡（model card）把预期用途、评价条件和已知限制写成模型报告，为这种使用判断提供依据\scite{mitchell2019modelcards}

方案相容以后，还要比较准备和持续使用各需要什么资源。表~\ref{tab:knowledge-acquisition-model-preparation}把标签建设、训练、接口改造与试验列为一次性投入，把推断资源、调用限制和版本维护列为持续条件。直接调用可能省去初始训练，却在每轮获取中引入数据传输、限额和服务版本变化；从头构建需要前期训练，通常也让我们有机会控制参数和部署范围。直接复用主要调整调用配置与输出约束，适配则还须确定允许改变哪些参数。这些是成本构成与控制范围，不是未经测量的方案优劣排序。

\begin{table*}[htbp]
  \centering
  \small
  \setlength{\RaggedRightParindent}{0pt}
  \renewcommand{\arraystretch}{1.15}
  \begin{tabularx}{\linewidth}{@{}>{\RaggedRight\arraybackslash}p{20mm}*{4}{>{\RaggedRight\arraybackslash}X}@{}}
    \toprule
    准备方案 & 起点与准备 & 一次性投入 & 持续条件 & 重点试验 \\
    \midrule
    从头构建 & 任务标签；训练并接入模型 & 标签、训练与调参 & 算力、维护与重训 & 标签覆盖、过拟合 \\
    直接复用 & 已有模型；对齐接口 & 接入、试标与合规检查 & 服务、限额与版本 & 语义错配、接口漂移 \\
    适配已有模型 & 预训练参数；确定更新范围 & 适配训练与能力检查 & 推断资源与再适配 & 负迁移、原能力下降 \\
    \bottomrule
  \end{tabularx}
  \caption{模型准备方案的资源与适用条件。按行比较起点、一次性投入、持续条件与待检验风险；三行描述准备动作，同一已有模型可由直接复用转入适配。接口漂移指服务版本或返回约定改变。负迁移指在同一当前任务中，用同一指标、相同评价样本与口径比较，使用迁移信息的方案比不使用该信息的方案表现更差；它与原模型能力下降是不同的检查对象。各成本仍需实际测量，表中不预设价格或效果名次。}
  \label{tab:knowledge-acquisition-model-preparation}
\end{table*}

把这些选择放回花卉任务：有少量物种标签时，可以试验轻量分类器，也可以利用已有视觉编码器，把照片转成特征后训练或适配分类部分；有同一物种体系的现成模型时，可以先直接试标；只有通用多模态模型时，则先检查它能否按任务说明、示范和输出约束使用当前类别表。三种条件下都要试验困难物种、低质量照片和类别表外对象，检查未覆盖类别、输入变化及高置信错误。这里的高置信错误指模型给出很高的接收得分，答案却与参照判断不符。选择阶段只得到带资源限制和假定范围的候选，是否采用还要经过试标。

\subsection{标注模型的训练与适配}
\label{sec:knowledge-acquisition-model-training-adaptation}

假设已有视觉表示可用，但输出层尚不能区分当前物种表。此时缺的是按任务答案调整参数的准则，而不是再加一段调用说明。我们用每张训练照片的预测与物种标签之间的损失衡量偏差，再在允许更新的参数配置中寻找平均偏差较小的配置。设训练集为$\mathcal{D}_{\mathrm{tr}}=\{(\bm{x}_i,y_i)\}_{i=1}^{n}$，$n>0$，$\bm{x}_i$为输入，$y_i$为任务答案，模型参数为$\bm{\theta}$，预训练起点为$\bm{\theta}_0$。若还希望限制更新对原参数的偏离或表达保留已有能力的偏好，可以在平均损失上加一项惩罚：
\begin{equation}
  \widehat{\bm{\theta}}
  \in
  \argmin_{\bm{\theta}\in\Theta_{\mathrm{upd}}}
  \frac{1}{n}\sum_{i=1}^{n}\ell\!\left(f_{\bm{\theta}}(\bm{x}_i),y_i\right)
  +\lambda\Omega(\bm{\theta},\bm{\theta}_0),
  \label{eq:knowledge-acquisition-adaptation-objective}
\end{equation}
其中$\Theta_{\mathrm{upd}}$是允许选择的参数配置集合，冻结的坐标在这个集合中保持原值；$f_{\bm{\theta}}$是相应参数下的预测模型，$\ell$衡量预测与答案的差异，$\Omega$表达保留偏好，$\lambda\geq0$控制惩罚强度。取训练样本的平均，是为了比较模型在这批任务数据上的总体拟合程度，不能据此保证每张照片都预测正确。从头训练可允许全部参数变化并去掉相对于$\bm{\theta}_0$的惩罚；只训练输出层时，表示部分不变；微调（fine-tuning）则继续调整部分或全部预训练参数。只有允许集合非空且目标能取到最小值时，式中的最小点集合才给出$\widehat{\bm{\theta}}$；实际迭代优化通常近似求解，仍须检查停止条件，不能由目标形式保证到达全局最小点。保留偏好也不等于原能力必然保住，需另在原能力对应的数据上检查。

这个目标能拟合什么，取决于放入训练集的对象与标签。若同一植株的连拍照片分散到训练、方案选择和最终检查三组，检查对象就与训练对象高度重复，结果不足以说明模型会识别新植株。因此，检验新植株时应将同一植株的全部照片归入同一划分组，再核对物种、拍摄条件和信息不足情形的覆盖。按采集批次划分只有在不同批次的植株互不重叠时才能实现植株隔离；若同株跨批次出现，就须将这些照片归入同组，不能仅凭批次不同认定是新植株。错误标签会把损失拉向错误答案，缺少某物种标签则没有提供学习该类别的任务信号；延长训练不能自动补出这些缺失。

数据分开后，还须分开三种判断。训练损失用于观察优化是否拟合当前目标；方案选择集用于比较冻结范围、权重和停止时刻；最终检查集只在这些决定冻结后使用。若某候选在方案选择集上的常见花卉正确率提高，低清照片上本应弃权却作答的比例也提高，就要同时记录这两种变化，不能只报平均正确率。选定版本应绑定训练数据、类别体系、参数更新范围和推断配置，后续记录才知道实际调用的是哪一组条件。

并非每个候选都需要上述优化。已有模型满足接口和试标要求时，冻结模型版本与调用配置即可，不另设参数训练。若模型根本不能接收所需图像或承载任务类别，少量微调也未必补得上缺口；这类问题应返回表~\ref{tab:knowledge-acquisition-model-preparation}重新选择方案，而不能仅凭训练损失下降继续使用。

\subsection{标注能力的试验与确认}
\label{sec:knowledge-acquisition-model-pilot-validation}

可调用版本还没有经过当前批次的能力检查。扩大使用前，我们先让它按正式流程处理一小批照片，这就是\term{模型试标}（model pilot annotation）：用小规模运行检查任务理解、输出解析、候选内容与处理成本。它与第~\ref{sec:knowledge-acquisition-human-pilot}节的人工试标都试运行完整任务，但作答主体和故障修正不同。样本要含典型情形、易混边界、低质量输入、类别表外对象和允许弃权的情形。每条记录保存样本标识、模型与输入版本、原始输出、解析状态、候选与参照答案；调用未完成、格式不合法和内容错误分别计数，才能区分接口故障、解析问题与判断问题。

根据这些分项结果，可以决定下一步。预定范围内的内容错误、弃权和调用成本满足预设批量准入要求时，采用当前方案；问题可由补标签、改参数或修接口处理时，继续适配；关键观测不可见或任务超出可确认范围时，补充人工、规则或外部文档等来源。试标集若反复参与模型、阈值和提示选择，就已是开发依据，不再提供方案冻结后的最终独立证明。准入指标、抽样及区间估计由第~\ref{chap:knowledge-quality-assessment-correction}章展开。

% 长作者引文从段中开始，需同时预留段首至标记和完整边注的高度。
\Needspace{100mm}
若试标后的候选还将用于训练另一个模型，就多了一层依赖：作答模型提供目标，受训模型按目标更新。前者在这里承担\term{任务教师}（task teacher）角色，也就是为学习者提供具有既定任务含义的目标；后者称为学生。名称不规定谁更大或更准确，学生已有预训练能力也不由教师标签完全决定。弱到强泛化研究按预训练计算规模区分较小的弱模型与较大的强模型：弱模型先用真值标签训练，再为留出实例产生标签，强学生用这些弱标签微调。在其自然语言处理、棋题和奖励建模实验中，强学生的同任务测试准确率可以高于弱监督者；但朴素微调的准确率仍低于同一强模型用真值标签微调的参照，差距随任务与模型规模而变\scite{burns2023weaktostrong}
这一研究观察不能直接移植到花卉任务。我们仍须在当前任务上分别检查教师与学生，不能推出学生总能纠正教师。

在花卉试标中，可同时放入山茶与茶梅、花瓣局部照片和类别表外植物，分别检查物种混淆、观测不足与范围识别。只在典型正面照片上获得可用结果，就只能确认这一范围；应为未知的照片被强行分类时，要补上弃权接口与原因，或增加其他来源。若准入使用置信度，还要检查分数与正确率的关系。\term{概率校准}（probability calibration）在这里考察：把接收分数解释成正确概率时，同一分数段的实际正确比例与平均预测概率有多大差异，而不是只比较分数大小。分段、参照答案和允许偏差要随评价范围确定，有限样本的正确比例还存在估计误差。校准关系既不保证单条高分答案正确，也不自动迁移到新地区或新模型版本。确认记录应保存可用范围、失败类型和人工接管条件，而不只保存总体分数。

\section{标注输入的组织与构造}
\label{sec:knowledge-acquisition-input-organization}

模型已经选定，下一步要确定它在本次作答中能看到什么。同一张花瓣照片，附有类别表与没有类别表、附有同株叶片与没有叶片，支持的判断并不相同。我们按输入的作用分开四类内容：任务要求规定问题与答案含义，样本及上下文提供当前观测，示范展示规范在其他实例中的用法，辅助证据补入外部事实或对象关系。四类功能可以组合，不要求每个模型接收全部内容；固定接口也可用任务配置、标签映射和外围检索组件实现所需功能。

图~\ref{fig:knowledge-acquisition-input-composition}中，两种线型把作答约定与当前判断依据分开。任务要求和示范主要说明怎么答，当前样本和辅助证据主要说明依据什么答。一张带答案照片可以同时展示字段格式并支持当前判断，此时要记录两种用途；否则，结果变化以后就无法辨认是格式引导改变了输出，还是新参考提供了物种信息。

下面为同一对象S042构造请求，依次加入任务、观测、示范和证据；把四段按顺序合并，即得到本例的完整输入。照片的可见内容写成观测文本，便于我们只凭书中文字核对哪些判断有依据。对象、名录、日期和输出都是教学构造，未调用真实模型，也不用于植物鉴定。

\begin{figure}[htbp]
  \centering
  % 图4-2：一次模型标注中的四类预测输入，按112 mm正文宽度设计。
\begin{tikzpicture}[
  x=1mm,
  y=1mm,
  font=\figurefont\footnotesize,
  text=BookInk,
  input block/.style={book node,draw=FigureInput,fill=FigureInput!6,line width=0.8pt,text width=37mm,minimum height=12mm,inner sep=1.4mm},
  model block/.style={book node,draw=FigureModel,fill=FigureModel!8,line width=0.9pt,text width=25mm,minimum height=18mm,inner sep=1.5mm},
  output block/.style={book node,draw=FigureOutput,fill=FigureOutput!6,line width=0.8pt,text width=17mm,minimum height=18mm,inner sep=1.2mm},
  convention/.style={knowledge arrow,draw=BookAmber,line width=0.8pt,dashed},
  evidence/.style={knowledge arrow,draw=BookBlue,line width=0.9pt},
  edge text/.style={fill=white,inner sep=0.7pt,font=\figurefont\scriptsize,text=BookMuted}
]
  \node[input block,draw=BookAmber] (requirement) at (21,0) {任务要求与预测输出约定\\Q1：类别C1、弃权、字段};
  \node[input block,draw=BookAmber] (demonstration) at (21,-18) {示范案例与参照答案\\S017：观测不足，弃权};
  \node[input block] (sample) at (21,-36) {待处理样本与必要上下文\\S042：花瓣局部，无叶片};
  \node[input block] (support) at (21,-54) {辅助证据与关联事实\\E03：名录R3的地块名录};

  \node[model block] (model) at (70,-27) {标注模型\\版本$M_5$};
  \node[output block] (output) at (101,-27) {S042\\样本标注\\未知、弃权\\证据与版本};

  \draw[convention] (requirement.east) -- node[edge text,above] {约定} ([yshift=6mm]model.west);
  \draw[convention] (demonstration.east) -- node[edge text,above] {示范} ([yshift=2mm]model.west);
  \draw[evidence] (sample.east) -- node[edge text,below] {当前观测} ([yshift=-2mm]model.west);
  \draw[evidence] (support.east) -- node[edge text,below] {外部证据} ([yshift=-6mm]model.west);
  \draw[knowledge arrow,draw=FigureModel,line width=1.0pt] (model.east) -- (output.west);
\end{tikzpicture}

  \caption{一次模型标注中的四类输入。虚线传入作答约定和示范，实线传入当前观测与外部证据；S042、O042与E03的标识限定各片段的对象和用途。四类内容按功能区分，不要求所有模型使用同一种提示接口。}
  \label{fig:knowledge-acquisition-input-composition}
\end{figure}

合并片段前，还须检查每段信息是否属于当前对象、是否占用过多输入空间、是否与其他片段冲突，以及是否暴露实际作答时不可获得的参照答案。检索负责找出候选材料，组织负责确定材料怎样对应当前对象并进入输入。没有区分这两步，模型用错证据时就难以判断是检索到了错材料，还是正确材料被关联到了错对象。

\subsection{任务要求与输出约定}
\label{sec:knowledge-acquisition-task-output-contract}

S042需要的是按指定物种表作出的判断，不是任意植物名称。任务要求因而先确定物种名录、杂交种与别名的处理、多张照片的对象对应，以及未知和不适用的条件。输出约定再确定怎样保存对象标识、内容、来源动作、原因和证据引用；能解释含义的模型得分另存。前者决定判断什么，后者让程序知道怎样读取结果。对能接受指令的模型，S042使用下列说明：

\begin{quote}
  请按规范Q1，仅根据本请求提供的观测和证据，判断S042所指植株的物种。任务范围限于类别表C1中的山茶、茶梅、月季。若作答，只返回其中一个规范名称，不自行补入别名、杂交种或其他类别。

  只有证据足以支持唯一类别时才作答，此时内容为物种名称，语义状态为“已给出”，动作为“作答”。观测不足、材料冲突或只能缩小候选范围时，内容留空，语义状态为“未知”，动作为“弃权”。只有已确认对象不属于本任务范围时，才记为“不适用”，内容和动作均留空。不得臆测不可见部位，也不得把某地出现过的物种直接当作该地每株植物的物种。

  只返回一个JSON对象，字段固定为object、content、semantic、action、reason、evidence，依次表示对象标识、内容、语义状态、来源动作、原因和证据编号列表。空值写为null，reason简述判断依据，不虚构概率分数。后附S017是完整示范，不是本次待作答对象，也不是S042的证据。
\end{quote}

这段请求让模型交付对象标识、候选内容及其语义状态、来源动作、理由和证据引用。执行是否完成由调用程序记录，指定属性是否核验由后续检查记录；模型自报成功不能代替运行日志，自报正确也不能代替核验。

固定分类接口不能直接执行这段文字，须把相同要求落实为配置。假设C1把索引$1\mapsto$山茶、$2\mapsto$茶梅、$3\mapsto$月季，模型给出向量$(0.55,0.40,0.05)$，最高分对应山茶；这个最大项只产生类别候选。再假设包装程序按最高分筛选，教学门槛设为$0.70$，本次最高分低于门槛，就返回未知并弃权，不把未知编码成物种索引。门槛在这里用来区分接收与弃权，是构造配置而非正确率保证。照片或特征仍按模型接口传入，示范和外部证据由支持它们的外围组件处理。规范变化时，单纯字段重命名可能允许迁移记录，类别拆分或含义改变则可能需要重新预测或复核。

规范还需要边界实例来检查含义是否一致。山茶与茶梅的正例、易混反例和未知例，可展示叶缘、花期只有在对应部位或时间已知时才参与判断。换成实体抽取时，输出不再是照片类别，而是原文片段及实体类型，还须约定它们怎样对应规范化名称。封闭分类只在给定类别集合中选择，开放文本允许集合外内容，多解任务则另须规定候选是否互斥。任务对象和输出空间改变后，“答案”字段的含义也必须随之说明。

\subsection{待处理样本与必要上下文}
\label{sec:knowledge-acquisition-sample-context}

类别要求确定以后，还须让模型知道判断的是哪株植物。待处理样本限定当前对象，上下文补入解释它所需、且实际使用时可获得的观测。照片要对应原图或稳定引用、判断区域、拍摄时间和对象标识；文本要对应原文位置与必要段落；语音和时序还要限定时间区间。先绑定对象再裁取上下文，才能把输出还原到正确位置，避免对邻近对象的判断写回当前样本。

现在向S042请求加入当前观测。用O042标识这一片段，是为了以后追查模型看过什么，不是提前提供物种标签：

\begin{quote}
  当前对象：S042，对应植株P042。观测O042：2026年4月12日在教学园地块P03拍摄的花朵局部，仅见粉红色花瓣；叶片、叶缘、完整花形与株形均不可见。地点和日期来自该照片的采集记录。本请求没有其他照片，也没有这株植物的已知物种标签。
\end{quote}

O042提供了地点和日期，没有提供叶缘。补入同一植株的叶片图会增加观测；只把提示改成“仔细看叶缘”，并不能让不可见的部位变得可见。

裁取方式决定哪些线索保留下来。花瓣局部可突出形态细节，却可能丢掉叶片和株形；整图补回这些线索，也会带入花盆或背景标牌等干扰。文本中相应的问题是切分边界：重叠窗口可保留跨界信息，但重复实体要按原文位置合并。时序任务还受判断时点约束，只能使用当时允许获得的历史；加入未来观测后，评价的就不再是相同条件下的在线预测。

合并多种观测时，对象、位置和时间必须一致。植株A的照片若配上植株B的地点，答案即使措辞合理也可能指向错对象；缺失信息若被默认值替代，记录也会误报模型所见的内容。因此要保存实际输入的各类信息。最终评价的上下文整理规则应在查看参照答案前冻结，否则按答案挑选“必要背景”，就可能把评价对象的答案泄漏进输入。

% 为首段及GPT-3完整作者边注共同预留空间。
\Needspace{125mm}
\subsection{示范案例与参照答案}
\label{sec:knowledge-acquisition-demonstrations}

S042已有作答规范和当前观测，但模型还可能不清楚证据不足时应怎样填写字段。可以加入另一对象已经完成的输入与参照答案，作为\term{示范案例}（demonstration example）：展示规则在实例上的用法，而不直接给出S042的物种。能采用\term{上下文学习}（in-context learning, ICL）的语言模型，可依据输入中的这些案例作答而不更新参数。GPT-3工作中的少样本设置把任务与示范放入文本交互，是这种用法的代表实例\scite{brown2020}
这里改变的是调用输入，不是模型权重；是否有效和输出是否稳定，仍需在当前模型上试验。

要展示的情况应与规范中的判断分支对应，字段也须与正式输入一致。只给清晰正面照片，没有说明观测不足时如何填写；加入未知示范能展示弃权分支，但不能保证模型就会正确弃权。下面为S042请求加入完整教学示范，答案按Q1的弃权条件构造：

\begin{quote}
  示范输入：对象S017，植株P017，类别表C1。观测O017：2026年4月10日在教学园地块P08拍摄，仅见失焦的粉红色花瓣局部，叶片和株形不可见；没有辅助证据。按同一规范Q1判断物种。示范输出如下：
\end{quote}
\begin{lstlisting}[style=book code,numbers=none]
{
  "object": "S017", "content": null,
  "semantic": "未知", "action": "弃权",
  "reason": "花瓣局部失焦，缺少区分物种的观测",
  "evidence": ["O017"]
}
\end{lstlisting}

S017这条记录展示证据不足时的内容、动作和原因写法。它不把S017与S042对应成同一植株，也不提供两株同物种的证据。示范理由同样要检查，因为错误理由会把错误的规范解释带进后续调用。

增加示范还会改变输入长度、排列和类别构成，不能只按条数判断作用。少样本提示的顺序敏感性研究在其分类实验中发现，同一组示范的不同排列会改变任务结果，一种模型上选出的顺序也未必适用于另一种模型\scite{knowledge-acquisition-prompt-order2022}
所以应保持任务与评价口径一致，在开发样本上比较候选示范并保存版本，不用最终评价样本反复挑选。示范占用输入空间和调用资源，比较时还须检查当前观测或证据是否被挤掉。

同一参考库也可用于另一种目的。固定案例若主要展示作答规则，就是本节的示范；按当前查询检索带标签照片、再让其答案参与计算，就是辅助证据，并可形成第~\ref{sec:knowledge-acquisition-neighbor-prediction}节的近邻预测。前者的选择依据是要展示哪些规范分支，后者是哪些参考支持当前对象；共享样本库不使两种用途等同。

\subsection{辅助证据与关联信息}
\label{sec:knowledge-acquisition-auxiliary-evidence}

示范说明了怎么答，S042仍缺少区分物种的当前依据。我们可以从文档、数据库、带标签参考样本或对象关系中补入\term{辅助证据}（auxiliary evidence），也就是当前观测之外、可能支持本次判断的材料。组织时要说明找什么、怎样与对象对应、取哪段内容，以及输出怎样回指原始来源。检索命中只表明查询与材料相似或字段匹配；材料是否支持一个物种答案，是下一项判断。

花卉任务可按照片特征寻找带标签参考，也可按地点与花期寻找植物名录。每条材料须保留来源、适用时间、对象标识和标签是否核验；同名植物可能不是同一对象，过时名录可能不再适用，转载也未必增加独立依据。文本实体同样需要用上下文消歧，即确定同一名称在当前材料中实际指谁，不能只按字符串相同合并事实。

对S042，我们固定检索字段为地块P03和日期2026年4月12日，取覆盖这两项条件的名录片段。把它接在示范后，就完成了四类输入的组合：

\begin{quote}
  辅助证据E03：教学园维护名录R3，第8条，记录时间为2026年4月1日。适用对象是地块P03，时间范围是2026年3月1日至6月30日。原文：“本期地块记录中有山茶和茶梅；此表不是完整物种清单，未建立逐株标识与物种的对应。”该片段已与R3核对，但没有核验P042的物种。请在这一范围内使用E03，只为当前对象S042返回一次结果。
\end{quote}

P042的地点与日期符合E03范围，所以该片段可以关联到当前请求。可是E03没有逐株对应，且明确不是完整清单，既不能确定P042物种，也不能排除清单外类别。与O042合起来，证据仍不足以支持唯一答案；按Q1，可得到下列构造输出：

\begin{lstlisting}[style=book code,numbers=none]
{
  "object": "S042", "content": null,
  "semantic": "未知", "action": "弃权",
  "reason": "O042缺少鉴别特征，E03没有逐株对应",
  "evidence": ["O042", "E03"]
}
\end{lstlisting}

若实际调用返回这条输出，程序将执行状态记为成功，来源记为模型$M_5$，取得条件保存Q1、C1、O042、S017和E03的引用，验证状态先记为候选。复核若确认弃权符合Q1，确认的属性是动作合理，物种内容仍为未知。再看另一条构造输出：它的JSON可以解析，却把证据说成了不同范围的事实：

\begin{lstlisting}[style=book code,numbers=none]
{
  "object": "S042", "content": "山茶",
  "semantic": "已给出", "action": "作答",
  "reason": "R3已经证明P042是山茶",
  "evidence": ["E03"]
}
\end{lstlisting}

这条答案把地块级名录改写为P042的逐株证明，应在内容核验中记为已拒绝，并保留原始输出。调用已经完成，不能因内容被拒绝而改记执行失败。前一条可作为符合规范的弃权记录，后一条不能作物种标签；决定差别的是证据能支持到哪一层，不是字段是否齐全。

这类材料也可接入检索增强生成（retrieval-augmented generation, RAG）：检索器先从外部索引取文档，生成模型再依据检索结果产生文本，把保存在参数外的文档记忆与参数模型结合起来\scite{knowledge-acquisition-rag2020}
它给出了一种检索与作答的组合，不定义所有辅助证据流程。材料怎样找到、所需证据是否找全、模型是否按材料范围作答，仍要分别检查；能生成答案不等于完成这些核验。

扩大检索范围可能找回原先遗漏的证据，也会带入无关内容、重复材料与额外调用成本。因此要约定候选数量，标明每份证据对应哪个样本，材料冲突时保留冲突而不任意覆盖。本节得到的是可追溯材料与关系，尚未规定怎样综合它们的答案。近邻和图结构的计算在第~\ref{sec:knowledge-acquisition-neighbor-prediction}、\ref{sec:knowledge-acquisition-graph-prediction}节展开，多来源结论的融合由第~\ref{chap:knowledge-fusion}章处理。

\section{候选知识的预测}
\label{sec:knowledge-acquisition-candidate-prediction}

输入已经绑定到当前对象，还需要说明它怎样变成答案。对同一张未标注照片，可以用训练好的分类器前向计算，可以找带标签的相似照片投票，也可以把它接入已知与未知照片组成的图，联立求出类别得分。图~\ref{fig:knowledge-acquisition-prediction-mechanisms}分别展示这三条路径：读取参数中的已学关系、读取近邻的参考答案、利用图上连接传播标签。学习表示、检索近邻与构图也可串在同一方案中；这里按答案主要怎样计算来区分机制，不按是否使用神经网络分类。

\begin{figure*}[htbp]
  \centering
  % 图4-3：三种预测机制对照，按169 mm全版心宽度设计。
\begin{tikzpicture}[
  x=1mm,
  y=1mm,
  font=\figurefont\footnotesize,
  text=BookInk,
  unknown/.style={circle,draw=FigureInput,fill=white,line width=0.9pt,minimum size=8mm,inner sep=0pt},
  known/.style={rectangle,draw=FigureInput,fill=FigureInput!10,line width=0.9pt,minimum size=8mm,inner sep=0pt},
  compute/.style={book node,draw=FigureModel,fill=FigureModel!8,line width=0.8pt,text width=16mm,minimum height=11mm,inner sep=1mm},
  answer/.style={book node,draw=FigureOutput,fill=FigureOutput!6,line width=0.8pt,text width=13mm,minimum height=11mm,inner sep=1mm},
  relation/.style={draw=BookMuted,line width=0.7pt},
  error relation/.style={draw=BookAmber,line width=0.8pt,dashed,rounded corners=1mm},
  calculation/.style={knowledge arrow,draw=FigureModel,line width=0.9pt},
  annotation/.style={font=\figurefont\scriptsize,text=BookMuted,align=center},
  edge label/.style={font=\figurefont\scriptsize,text=BookMuted,fill=white,inner sep=.5pt}
]
  \draw[BookRule,line width=0.7pt] (55,2) -- (55,-52);
  \draw[BookRule,line width=0.7pt] (110,2) -- (110,-52);

  \node[anchor=north west,font=\figurefont\bfseries] at (1,1) {A\quad 已学映射};
  \node[unknown] (direct-query) at (9,-25) {S};
  \node[annotation] at (9,-34) {查询S042};
  \node[compute] (direct-model) at (29,-25) {预测模型\\$f_{\bm{\theta}}$};
  \node[answer] (direct-answer) at (48,-25) {候选\\山茶};
  \draw[calculation] (direct-query.east) -- (direct-model.west);
  \draw[calculation] (direct-model.east) -- (direct-answer.west);
  \node[annotation] at (29,-44) {预测时主要读取模型参数};

  \node[anchor=north west,font=\figurefont\bfseries] at (57,1) {B\quad 近邻参考};
  \node[unknown] (neighbor-query) at (64,-27) {S};
  \node[annotation] at (64,-37) {查询S042};
  \node[known] (neighbor-a1) at (79,-13) {A};
  \node[known] (neighbor-a2) at (79,-27) {A};
  \node[known] (neighbor-b) at (79,-41) {B};
  \node[compute,text width=13mm] (vote) at (98,-27) {加权\\投票};
  \node[answer,text width=13mm] (neighbor-answer) at (98,-46) {候选\\山茶};
  \draw[relation] (neighbor-query) -- (neighbor-a1);
  \draw[relation] (neighbor-query) -- (neighbor-a2);
  \draw[relation] (neighbor-query) -- (neighbor-b);
  \draw[calculation] (neighbor-a1.east) -- (vote.west);
  \draw[calculation] (neighbor-a2.east) -- (vote.west);
  \draw[calculation] (neighbor-b.east) -- (vote.west);
  \draw[calculation] (vote.south) -- (neighbor-answer.north);
  \node[annotation] at (79,-51) {$A=$山茶，$B=$茶梅};

  \node[anchor=north west,font=\figurefont\bfseries] at (112,1) {C\quad 图结构};
  \node[known] (graph-a) at (119,-11) {$A:1$};
  \node[known] (graph-b) at (119,-42) {$B:0$};
  \node[unknown] (graph-u) at (137,-27) {$u$};
  \node[unknown] (graph-v) at (151,-11) {$v$};
  \node[answer,text width=14mm] (graph-answer) at (160,-33) {$u$的A分数\\$\frac23\to\frac{11}{15}$\\错边后$\frac{17}{27}$};
  \draw[relation] (graph-a) -- node[edge label,above] {$2$} (graph-u);
  \draw[relation] (graph-b) -- node[edge label,below] {$1$} (graph-u);
  \draw[relation] (graph-u) -- node[edge label,above] {$1$} (graph-v);
  \draw[relation] (graph-a) -- node[edge label,above] {$3$} (graph-v);
  \draw[error relation] (graph-b.east) -- (graph-v.south |- graph-b.east) -- (graph-v.south);
  \node[annotation,text=BookAmber] at (137,-45) {错误边$+3$};
  \draw[calculation] (graph-u.east) -- (graph-answer.west);
  \node[annotation] at (136,-51) {方形为标签锚点；边权可按例~\ref{ex:knowledge-acquisition-graph-propagation}复算};
\end{tikzpicture}

  \caption{三种预测机制怎样使用已有知识。圆形是未知答案的查询或节点，方形是带答案参考或标签锚点；无箭头边表示相似关系，箭头表示计算路径。面板A、B示意候选产生，面板C对照直接邻居、多步传播和错误连边的得分。类别、二维位置、距离与边权均为教学构造，未应用Q1接收约定，不改写S042的弃权判断。}
  \label{fig:knowledge-acquisition-prediction-mechanisms}
\end{figure*}

三条路径需要保留的依据不同，表~\ref{tab:knowledge-acquisition-prediction-mechanisms}将这些信息与核心计算、失效条件对应起来。直接模型可先产生表示，近邻再提供参考，图推断再协调一批节点，三者不构成互斥模型族。采用组合方案时，中间的参数映射、参考选择和图连接都应可检查；相同的最终类别不能说明中间依赖也相同。

\begin{table*}[htbp]
  \centering
  \small
  \setlength{\RaggedRightParindent}{0pt}
  \renewcommand{\arraystretch}{1.15}
  \begin{tabularx}{\linewidth}{@{}>{\RaggedRight\arraybackslash}p{20mm}*{4}{>{\RaggedRight\arraybackslash}X}@{}}
    \toprule
    预测机制 & 读取的信息 & 核心计算 & 新样本处理 & 主要依赖与风险 \\
    \midrule
    已学映射 & 参数、接口与当前输入 & 前向计算或解码 & 按接口推断 & 训练覆盖、接口与分布变化 \\
    近邻参考 & 参考答案、表示与索引 & 距离排序与加权 & 检索并聚合 & 距离含义、标签质量与检索遗漏 \\
    图结构 & 边权、标签锚点与节点 & 联立求解或传播 & 接入图并更新解 & 错边、错标签与无锚点分量 \\
    \bottomrule
  \end{tabularx}
  \caption{三种预测机制的使用条件。各行按当前答案的主要计算路径，列出读取的信息、新样本的处理及依赖风险；学习表示、检索与图推断可在同一方案中组合，不是模型结构的永久分类。}
  \label{tab:knowledge-acquisition-prediction-mechanisms}
\end{table*}

\subsection{基于已学映射的直接预测}
\label{sec:knowledge-acquisition-direct-prediction}

先看一个已训练的花卉分类器。它不必在每次调用时重新找带标签照片，而把当前输入通过已固定的参数映射到类别得分；本章称这种路径为\term{直接预测}（direct prediction）。设$\mathsf{u}_i$是样本$i$及接口允许的任务要求、上下文和辅助内容，$\bm{\theta}$为调用版本的参数。类别表规定候选集合$Y=\{1,\ldots,K\}$，$s_{\bm{\theta},c}(\mathsf{u}_i)$是类别$c$的得分，约定分数越高表示模型越支持该类别。要从这些得分得到一个类别候选，可以选取得分最大的类别：
\begin{equation}
  \widehat{y}_i
  \in
  \argmax_{c\in Y}s_{\bm{\theta},c}(\mathsf{u}_i)
  \label{eq:knowledge-acquisition-direct-classification}
\end{equation}
最大值只完成同一类别表内的排序；若并列，要按预定规则选取或保留多个候选。需要弃权时，还须规定最高分、最高与次高分的差值等接收条件。分数只有经模型定义与相应评价支持，才能解释成概率；即使校准关系成立，$\argmax$也不证明这张照片的答案正确。

同一参数映射还可以在另一种监督单位下训练。第~\ref{sec:rule-based-knowledge-acquisition-relation-matching}节的远程监督给出同一实体对的提及组，而没有确定哪句表达关系。沿用两句话：句乙“林舟在论坛上批评了海岬科技的发布流程”只是共同提及；句甲“海岬科技任命林舟为首席架构师”表达任职。语义对照是教学解释，训练时只给“组内至少一句表达任职”的组级监督，不给这两个逐句答案。因此，句级预测还需要一个组级聚合量，才能与现有监督比较。

令$B$为非空提及组的句子索引，$Z_j\in\{0,1\}$表示句$j$是否表达目标关系。模型从句子和实体标识组成的$\mathsf{u}_j$计算$p_j=p_{\bm{\theta}}(Z_j=1\mid\mathsf{u}_j)$。我们需要“至少一句表达”的概率，可以先算其补事件“一句也不表达”。\term{噪声或}（noisy-or）正是这样聚合：假定给定整组输入后各$Z_j$条件独立，且各自概率只依赖对应$\mathsf{u}_j$，于是补事件概率为各$(1-p_j)$之积。再用组标签$y_B\in\{0,1\}$与组级概率比较，二元交叉熵惩罚模型给参照事件的低概率；采用自然对数时，
\begin{equation}
  \begin{aligned}
    \widehat p_B
    &=1-\prod_{j\in B}(1-p_j),\\
    \ell_B
    &=-y_B\log\widehat p_B
      -(1-y_B)\log(1-\widehat p_B).
  \end{aligned}
  \label{eq:knowledge-acquisition-model-bag-prediction}
\end{equation}
上述对数式在$0<\widehat p_B<1$时可直接计算。端点按极限解释，零系数的对数项取0：参照事件的预测概率为1时损失为0，为0时损失为$+\infty$。若数值实现将概率截断到$(0,1)$内以避免无穷值，须另记截断规则，不能把截断后的有限值当作原端点损失。
例如构造一个正组，按句乙、句甲的顺序设$p_1=0.1$、$p_2=0.8$、$y_B=1$。两句都不表达的概率为$0.9\times0.2$，因此
\[
  \widehat p_B=1-(1-0.1)(1-0.8)=0.82,
  \qquad
  \ell_B=-\log0.82\approx0.19845.
\]
训练梯度通过聚合式传给句级模型，拟合的是整组的“至少一句”事件，不是要求每句都为正。预测可返回组分数$0.82$，并把得分更高的句甲列为证据候选；$0.82$不是句甲或句乙各自的概率。多个句级概率配置可以形成相同组分数，因而仅有组标签不能唯一找出支撑句。

补事件相乘依赖刚才的条件独立假设。重复转载或高度相似句子可能共享同一证据，按独立句相乘就可能高估组级支持。另一个问题在监督来源：知识库关系存在，不保证当前组有句子表达它。若$y_B=1$仍是远程监督候选，应先决定接收、降权或复核；上式只展示接受该监督之后怎样训练，没有把它核验为事实。

从固定类别或组级事件转到其他任务，参数映射的输出含义也随之改变。回归给出数值，结构化预测还要满足序列、树或集合的组合约束。若让模型生成物种名称和说明，候选不再只是$Y$中的一个索引，而是词元表中词元组成的序列。设输入仍为$\mathsf{u}_i$，输出长度为$m$，第$r$个词元为$y_r$。要计算完整序列的条件概率，可按概率的链式法则，将它分解为给定输入和已生成前缀的逐步概率：
\begin{equation}
  p_{\bm{\theta}}(y_{1:m}\mid\mathsf{u}_i)
  =
  \prod_{r=1}^{m}
  p_{\bm{\theta}}(y_r\mid y_{1:r-1},\mathsf{u}_i).
  \label{eq:knowledge-acquisition-direct-generation}
\end{equation}
这个乘积说明每一步如何依赖前缀，本身还没有选定最终文本。解码（decoding）是按这些条件概率选取或采样词元以形成序列的过程，策略、候选数量、终止条件与最大长度都会影响结果。生成的物种名还要映射到类别表并处理别名。通顺说明只表明模型完成了文本生成；证据是否真实、是否支持该物种，仍要另查。

序列或类别生成以后，批量流程还要解析输出并还原对象。重叠文本窗口可能给同一实体不同答案，损坏图像可能返回空结果，字段也可能语法合法而含义错误。保存原始输出与解析候选，可以分别检查模型输出、解析和对象写回。增加采样或候选数量会扩大待处理集合，也增加融合与复核工作；多份候选的综合由第~\ref{chap:knowledge-fusion}章处理。

还可以改变生成对象：不固定一批待标照片，而让模型共同提出任务、输入和答案。\term{Self-Instruct}采用这种材料构造方式，以已有种子任务生成新指令和输入输出实例，再过滤无效或相似材料，用保留实例继续微调\scite{wang2023selfinstruct}
此时模型不仅改变标签，还改变训练任务与输入的分布。检查因而要覆盖任务是否成立、输入是否符合指令、答案是否可用和材料是否重复；固定样本标签检查不能全部沿用。生成更多实例只增加候选数，是否增加任务覆盖或训练效用需另有证据。

\subsection{基于近邻参考的预测}
\label{sec:knowledge-acquisition-neighbor-prediction}

假设当前花卉分类器的输出还不适用，但我们保存了一批带物种标签的照片，可以比较新照片与它们的形态。此时答案可从调用时找到的参考项取得，而不只从参数读取。\term{近邻参考预测}（neighbor-reference prediction）就是选择与查询相近的带答案样本，再聚合其答案。设参考集为$\mathcal{R}=\{(\bm{x}_j,y_j)\}_{j=1}^{n}$，$y_j$属于任务类别表$Y$，$\phi$将照片映射到比较空间，$d$是该空间的距离。按距离从小到大选出$k$个参考索引$N_k(\bm{x})$，其中$1\leq k\leq n$。要让参考标签形成可比较的类别支持，可先给邻居非负权重并归一化，再选取总支持最大的类别：
\begin{equation}
  \begin{aligned}
    \widehat{y}(\bm{x})
    &\in
    \argmax_{c\in Y}
    \sum_{j\in N_k(\bm{x})}
    \alpha_j(\bm{x})\mathbf{1}\{y_j=c\}, \\
    \alpha_j(\bm{x})
    &=
    \frac{w\!\left(d(\phi(\bm{x}),\phi(\bm{x}_j))\right)}
    {\sum_{r\in N_k(\bm{x})}w\!\left(d(\phi(\bm{x}),\phi(\bm{x}_r))\right)}.
  \end{aligned}
  \label{eq:knowledge-acquisition-neighbor-prediction}
\end{equation}
其中$w$把距离转成非负权重，通常令距离越小权重越大，也可取等权；分母要求至少一个邻居有正权重。归一化让权重之和为1，每类得分就是分给该类的参考权重，而不是自动校准的正确概率。回归时参考答案$y_j$改为数值，保留邻域与归一化权重，直接取$\widehat y(\bm{x})=\sum_{j\in N_k(\bm{x})}\alpha_j(\bm{x})y_j$作为加权平均预测，不再对类别集合取$\argmax$。距离并列、类别平票、无有效邻居和查询与参考完全重复都须先规定处理；与最终评价对象重复的参考还要排除或明确评估口径，不能把查回自身标签当作泛化。

经典近邻分类以局部参考形成预测\scite{coverhart1967}
本节已给出检索与聚合。即使$\phi$由神经模型学习，式~\eqref{eq:knowledge-acquisition-neighbor-prediction}仍读取参考标签。更正或撤回一个参考标签可能直接改变附近查询，不要求先重训表示模型；这与只调用固定参数的映射不同。近邻算法与统计性质可参考\BookChapterRef{chap:classification}{卷四《模型：经典模型》的“分类”章}。

这条路径能提供物种依据，依赖距离确实反映任务相关特征。若表示主要编码背景与设备，公园里的山茶可能检索到同地点的月季；只有表示保留足够鉴别特征时，形态近邻才可能支持物种判断。标签一致还可能来自同一错误来源，类别失衡也会使多数类挤占邻域。参考稀疏时，“这批中最近”不等于“足够相似”；应按预定距离与覆盖要求降低接收支持或弃权。

保留更多参考也带来存储、索引构建与查询成本。近似近邻索引（approximate nearest-neighbor index）用可能遗漏最近项的搜索换取较少计算，遗漏会改变投票；新增或撤回参考还可能要求维护索引与缓存。候选记录因此应保存实际采用的参考标识、距离、权重和索引版本，后续才能重放式~\eqref{eq:knowledge-acquisition-neighbor-prediction}，区分参考答案变更与检索实现变化。

\subsection{基于图结构的预测}
\label{sec:knowledge-acquisition-graph-prediction}

若一批未标注照片互相相似，只汇总每张查询附近的已知标签，会忽略未知照片之间的连接。我们可以将已知与未知样本共同放入图，让连接约束多个节点的答案，这就是\term{图结构预测}（graph-based prediction）。这里采用的假设是：任务相关的强连接应对应相近的类别得分，并由已知标签确定类别含义。\term{标签传播}（label propagation）沿这些连接传递标签信息，不把相似边解释成因果关系。设无向图$\mathcal{G}=(V,E)$有非负对称边权$w_{i,j}=w_{j,i}$，无边处权重为0；$V$分为已标节点$L$与未知节点$U$。已标节点称为标签锚点，因为它们提供类别参照。用$\bm{F}_i\in\mathbb{R}^{K}$保存节点$i$的$K$类得分，所有行组成$\bm{F}\in\mathbb{R}^{|V|\times K}$，已知标签$\bm{Y}_i$为独热向量，即对应类别为1、其余为0。我们需要一组未知得分，既保持$\bm{F}_L=\bm{Y}_L$，又让强连接两端差异小，因此在满足边界的得分矩阵中最小化加权平方差：
\begin{equation}
  \min_{\bm{F}:\,\bm{F}_L=\bm{Y}_L}
  \frac{1}{2}\sum_{i,j\in V}w_{i,j}\lVert\bm{F}_i-\bm{F}_j\rVert_2^2
  =
  \min_{\bm{F}:\,\bm{F}_L=\bm{Y}_L}
  \operatorname{tr}\!\left(\bm{F}^{\top}\bm{L}_{\mathcal{G}}\bm{F}\right),
  \label{eq:knowledge-acquisition-graph-energy}
\end{equation}
其中$\bm{W}$是边权矩阵，$\bm{D}$是对角矩阵，$\bm{D}_{i,i}=\sum_jw_{i,j}$；图Laplacian矩阵$\bm{L}_{\mathcal{G}}=\bm{D}-\bm{W}$把同一平方差写成矩阵形式，$\operatorname{tr}$表示对角元素之和。对称权重使$(i,j)$和$(j,i)$重复计数，$1/2$抵消重复；平方差展开后的两项节点平方范数合并为度数项，内积项由$\bm{W}$给出，便得到右侧。最小化的是图上加权不一致，不是未知样本的真实分类错误；其作用依赖连边符合刚才的任务相似假设。

要从目标得到实际未知得分，按$L,U$排列行列并固定$\bm{F}_L=\bm{Y}_L$。未知部分的梯度为$2(\bm{L}_{U,U}\bm{F}_U+\bm{L}_{U,L}\bm{Y}_L)$，令其为0，便得到线性方程；可逆时再写成逆矩阵形式：
\begin{equation}
  \bm{L}_{U,U}\bm{F}_U
  =
  -\bm{L}_{U,L}\bm{Y}_L,
  \qquad
  \bm{F}_U
  =
  -\bm{L}_{U,U}^{-1}\bm{L}_{U,L}\bm{Y}_L.
  \label{eq:knowledge-acquisition-graph-solution}
\end{equation}
可逆性不是对任意图的保证。对非负无向权重，若每个含未知节点的正权重连通分量都能到达标签锚点，固定边界后未知部分有唯一解。其依据是：保持标签边界不变时，若未知得分的变化具有零能量，就须在每个正权重连通分量内为常数；这项变化在锚点为0，故只能全为0，未知部分的二次型因而正定。完全没有锚点的未知分量可任取常数得分，能量仍为0，绝对类别无法确定，相应矩阵也奇异。Gaussian fields与调和函数方法给出了这种构造、矩阵解及随机游走解释\scite{zhu2003}

同一线性方程还给出局部计算方式。把未知节点$i$的驻点条件移项，得到它的得分等于邻居得分的加权平均，这称为调和条件。采用同步边界迭代时，可以将未知节点初始化为均匀得分$\bm{F}_i^{(0)}=(1/K,\ldots,1/K)$，已知节点始终固定为$\bm{Y}_i$；每轮所有未知节点都读取上一轮的邻居得分，同时计算新值：
\[
  \bm{F}_i^{(t+1)}
  =
  \frac{\sum_j w_{i,j}\bm{F}_j^{(t)}}{\sum_j w_{i,j}},
  \qquad i\in U,
\]
其中分母要求节点有正权重邻居。在上述锚点条件下，从任意有限初值出发，这种同步迭代都收敛到唯一调和解：未知节点上的初值影响沿连接逐步传向固定锚点，不会永久困在无锚点分量中。实际计算须预先规定近似停止标准，例如当所有未知节点的相邻两轮得分变化的最大二范数不超过容差$\varepsilon>0$时停止，并将最大轮数$T_{\max}$设为正整数。先达到最大轮数而变化仍超容差时，应另记求解未收敛，不把当前得分当作已取得调和解；也可直接求解前述线性方程。最后仍需从每行类别得分选出最高项，并应用预定平票与弃权规则，得分矩阵本身还不是已核验的类别标签。与只聚合查询附近已知标签相比，未知节点也能把邻居信息传到后续节点。范围扩大，错误锚点和跨类别错边的影响也可能传得更远；数值收敛只说明求解到图目标的解，不说明图假设正确。

\begin{example}[直接近邻、多步传播与错误连边]
  \label{ex:knowledge-acquisition-graph-propagation}
  用一个构造图检查未知节点的连接改变了什么，只计算类别A的一维得分。锚点A、B分别取$f_A=1$、$f_B=0$；查询$u$以权重2、1连接A、B，以权重1连接未知$v$；$v$再以权重3连接A。若只读取$u$直接相连的已标答案，忽略$v$并归一化，得到
  \[
    f_u^{\mathrm{direct}}=\frac{2f_A+f_B}{2+1}=\frac{2}{3}.
  \]
  完整图则保留$u$与$v$的连接。两节点的总边权都为4，调和条件给出
  \[
    4f_u=2+f_v,\qquad 4f_v=3+f_u,
  \]
  联立得$f_u=11/15$、$f_v=14/15$，说明A锚点经$v$影响了$u$。再人为加入错误边：让$v$以权重3连接B，第二式变为$7f_v=3+f_u$，解变为$f_u=17/27$、$f_v=14/27$。假设教学规则仅当A得分至少$0.7$时接收A候选，原图中的$11/15$达到阈值，错边后的$17/27$没有达到。阈值在此展示得分变化如何影响接收，不是实际质量标准。图~\ref{fig:knowledge-acquisition-prediction-mechanisms}的面板C使用同一组节点、权重和结果。
\end{example}

花卉图中，照片是节点，形态表示、地点或同株关系可以提供构边线索，但每类线索仍须检查是否支持物种相似。背景相似而物种不同的照片被强连接时，错误可能沿图扩大；加入新照片改变邻接关系，也可能改变旧得分。构图、稀疏存储和求解都有成本，实际可解线性方程而不显式求逆。即使表示与边权可学习，本节仍由已有物种标签给输出赋予类别含义。第~\ref{chap:data-derived-knowledge-acquisition}章的图分组则先发现未赋予物种语义的群组；同样使用图，不等于取得同一种知识。

\section{获取过程的迭代改进}
\label{sec:knowledge-acquisition-iterative-improvement}

一轮花卉候选已经形成，复核可能发现类别边界学错、叶片观测缺失、名录关联错误，也可能发现预算大多用于重复处理简单照片。前三者的修改对象分别在模型或输入中，最后一项则涉及下一轮处理哪些样本。迭代前先保存并冻结上一轮的数据、模型、输入、索引或图版本与输出，再按成因选择更新模型、输入或获取安排。修改以后，还要沿版本依赖判断哪些旧记录需重算，而不能默认重训一次就处理了所有问题。

选择修改方式所用的数据，是开发验证和过程审计依据：它们帮助确定阈值、输入版本、教师更新方法及继续、停止或回退。参与这些选择后，同一批数据不能再充当最终独立验收。图~\ref{fig:knowledge-acquisition-iterative-update}上方面板将版本决策与留出验收分开，下方面板展示一个教师更新过程。训练目标下降只说明当前训练信号拟合得更好，所以不能直接连到候选质量改善。

\subsection{模型的更新}
\label{sec:knowledge-acquisition-model-update}

假设已取得少量人工物种标签，模型又为未标注照片提出了候选，我们希望让其中可接收的一部分参与下一轮训练。\term{自训练}（self-training）就是先用当前模型预测，再把筛选后的预测暂时当作目标来更新模型。这类临时目标称为\term{伪标签}（pseudo-label），它们来自模型候选，不因进入训练就成为真实标签。设人工标注索引为$L$，第$t$轮未标注候选索引为$U^{(t)}$，二者分开保存；模型给出类别$\widehat y_j^{(t)}$及接收分数$q_j^{(t)}$，约定分数越高越符合接收偏好。要明确本轮哪些候选进入训练，可以用在开发阶段选择的阈值$\tau^{(t)}$限定集合：
\begin{equation}
  P^{(t)}
  =
  \left\{
    \left(\bm{x}_j,\widetilde y_j^{(t)},q_j^{(t)}\right)
    :
    j\in U^{(t)},\
    q_j^{(t)}\geq\tau^{(t)},\
    \widetilde y_j^{(t)}=\widehat y_j^{(t)}
  \right\}.
  \label{eq:knowledge-acquisition-pseudo-label-set}
\end{equation}
集合中的$\widetilde y_j^{(t)}$取当前预测类别，并连同样本与接收分数保存。阈值回答哪些候选被采用，不回答每条是否正确；分数若未作相应概率校准，只能按当前模型的尺度解释筛选。即使总体校准成立，也须检查所接收类别和条件的覆盖，不能把高分筛选等同于独立核验。

接收集合确定后，还缺少它与人工标签共同影响参数的目标。对$|L|>0$且$P^{(t)}$非空的情形，可分别计算人工标签损失与伪标签加权损失，再用系数控制两类信号的相对影响。这里选择样本平均而非损失总和，是为了不让接收数量本身直接决定整批信号的量级。一种示意目标为
\begin{equation}
  \mathcal{L}^{(t)}(\bm{\theta})
  =
  \frac{1}{|L|}\sum_{i\in L}\ell\!\left(f_{\bm{\theta}}(\bm{x}_i),y_i\right)
  +
  \lambda^{(t)}
  \frac{1}{|P^{(t)}|}
  \sum_{(\bm{x}_j,\widetilde y_j^{(t)},q_j^{(t)})\in P^{(t)}}
  a_j^{(t)}
  \ell\!\left(f_{\bm{\theta}}(\bm{x}_j),\widetilde y_j^{(t)}\right),
  \label{eq:knowledge-acquisition-self-training-objective}
\end{equation}
其中$a_j^{(t)}\geq0$控制单条伪标签的影响，$\lambda^{(t)}\geq0$控制伪标签项相对人工项的系数。按固定的可更新参数范围优化这个目标后，得到下一轮模型；该目标惩罚的是对当前训练答案的偏离，错误伪标签也会被拟合。两类损失分别按各自样本数归一化，避免统一平均时仅因伪标签变多就改变两类数量占比，实际量级仍受$a_j^{(t)}$与损失大小影响。$P^{(t)}$为空时跳过伪标签项，并据开发证据停止、调整接收规则或补来源，不用除零约定掩盖没有接收样本的事实。

% 本段末尾两条边注不能跨页续排，需连同段落检查可用高度。
\Needspace{110mm}
把接收与训练连起来，就得到“预测候选—选择或复核—更新参数—重新预测”的循环。每轮保存人工标签、伪标签、权重与模型版本，并说明历史伪标签固定、重估还是允许撤销。一次加入很多标签可扩大训练覆盖，也可能同时放大一类错误；分批加入便于检查轮间变化，代价是更多预测和训练。继续训练沿用当前参数与优化状态，重新训练可从指定起点重走训练路径，二者都须按当前任务比较。早期自训练思想已用于自适应模式识别\scite{scudder1965selftraining}
现代深度模型研究还观察到确认偏差（confirmation bias）：模型的错误伪标签被当作训练答案，更新又强化该错误\scite{arazo2020confirmation}

循环中还须确定下一轮由哪个模型提供任务目标，这就是任务教师的更新安排。固定教师始终提供同一版本预测，便于追溯但不吸收学生变化；定期替换在检查后把某个学生版本升为教师；随学生更新的教师则更频繁地调整参数。定期替换允许学生使用更大或不同结构，但角色互换不证明预测更准确，新教师也可能继续提供旧错误。

如果每一步都直接把学生复制为教师，目标也会跟随学生的当前参数变化。\term{Mean Teacher}改为保留一个加权汇总历史学生参数的教师，并让同一样本的学生预测与教师预测保持一致\scite{knowledge-acquisition-mean-teacher2017} 其中指数移动平均（exponential moving average, EMA）每步保留一部分旧教师参数，再加入一部分新学生参数；它改变目标更新速度，不提供目标正确性的保证。

为了算出这一更新，需要同时定义预测差异与参数变化。第$t$步用同一样本的两种允许扰动$\widetilde{\bm{x}}_s,\widetilde{\bm{x}}_t$分别输入学生和教师；扰动须按任务设计，不应改变期望保持的标签含义。差异函数$d$取值越小表示两项预测越一致，加上真实标签的监督损失，得到学生要减小的目标。学生参数为$\bm{\theta}^{(t)}$，同结构教师参数为$\overline{\bm{\theta}}^{(t)}$，初始化可从同一模型复制；当前梯度只更新学生，然后才按EMA更新下一步教师：
\begin{equation}
  \begin{aligned}
    \mathcal{L}_{\mathrm{con}}^{(t)}
      \bigl(\bm{\theta},\overline{\bm{\theta}}^{(t)}\bigr)
    &=
    d\!\left(
      f_{\bm{\theta}}(\widetilde{\bm{x}}_s),
      f_{\overline{\bm{\theta}}^{(t)}}(\widetilde{\bm{x}}_t)
    \right), \\
    \mathcal{J}^{(t)}
      \bigl(\bm{\theta};\overline{\bm{\theta}}^{(t)}\bigr)
    &=
    \mathcal{L}_{\mathrm{sup}}^{(t)}(\bm{\theta})
    +\lambda_{\mathrm{con}}^{(t)}
      \mathcal{L}_{\mathrm{con}}^{(t)}
      \bigl(\bm{\theta},\overline{\bm{\theta}}^{(t)}\bigr), \\
    \bm{\theta}^{(t+1)}
    &=
    \bm{\theta}^{(t)}
    -\eta^{(t)}
      \left.
      \nabla_{\bm{\theta}}\mathcal{J}^{(t)}
      \bigl(\bm{\theta};\overline{\bm{\theta}}^{(t)}\bigr)
      \right|_{\bm{\theta}=\bm{\theta}^{(t)}}, \\
    \overline{\bm{\theta}}^{(t+1)}
    &=
    \alpha\overline{\bm{\theta}}^{(t)}
    +(1-\alpha)\bm{\theta}^{(t+1)}.
  \end{aligned}
  \label{eq:knowledge-acquisition-mean-teacher}
\end{equation}
其中$\mathcal{L}_{\mathrm{sup}}^{(t)}$是当前真实标签的监督损失，$\lambda_{\mathrm{con}}^{(t)}\geq0$平衡一致性项，$\eta^{(t)}>0$为学生步长，$\alpha\in[0,1)$决定保留旧教师的比例。选择一致性损失，是要约束同一样本两次任务预测的差异；真实标签项仍提供外部任务参照。教师在求梯度时固定，最后一行才改变它，下一步由$f_{\overline{\bm{\theta}}^{(t+1)}}$重算目标。移动平均作用于参数，不是每个样本的历史预测；一致性下降也可能只是学生更接近错误教师，仍需独立检查。

我们用单个标量参数检查最后一行。假设$\alpha=0.9$、旧教师参数为$0.6$、梯度更新后的学生参数为$1.0$，新教师就是$0.9\times0.6+0.1\times1.0=0.64$。再把输入交给$f_{0.64}$才得到下一轮目标。模型通常对参数非线性，故这次参数平均一般不等于预测平均$0.9f_{0.6}+0.1f_{1.0}$。

上例只算出了新参数，尚没有决定接收哪些硬标签。若把教师再接入式~\eqref{eq:knowledge-acquisition-pseudo-label-set}，还要重算类别与接收分数；Mean Teacher原有的软一致性目标不要求这一步硬阈值筛选。为展示组合，另行假设同一候选更新前由$f_{0.6}$给出$0.72$，更新后由$f_{0.64}$给出$0.78$，阈值都为$0.75$。于是前一轮不接收，下一轮才把教师类别作为硬伪标签放入$P^{(t+1)}$。分数是额外构造值，不能由$0.6$、$0.64$推出，更不说明分数随参数单调增加。

\begin{figure*}[!t]
  \centering
  % 图4-4：上下两面板共享169 mm版心；模型参数读取、梯度与EMA使用独立路由。
\begin{tikzpicture}[
  x=1mm,
  y=1mm,
  font=\figurefont\footnotesize,
  text=BookInk,
  loop node/.style={book node,draw=FigureModel,fill=FigureModel!7,line width=0.8pt,text width=24mm,minimum height=12mm,inner sep=1.2mm},
  data node/.style={book node,draw=FigureInput,fill=FigureInput!7,line width=0.8pt,text width=24mm,minimum height=12mm,inner sep=1.2mm},
  audit node/.style={book node,draw=BookAmber,fill=BookAmber!7,line width=0.8pt,text width=24mm,minimum height=12mm,inner sep=1.2mm},
  state node/.style={book node,draw=BookMuted,fill=BookPaper,line width=0.7pt,text width=24mm,minimum height=12mm,inner sep=1.2mm},
  loss node/.style={book node,draw=FigureLoss,fill=FigureLoss!7,line width=0.8pt,text width=24mm,minimum height=12mm,inner sep=1.2mm},
  main flow/.style={knowledge arrow,draw=FigureModel,line width=0.9pt,rounded corners=1mm},
  evidence flow/.style={knowledge arrow,draw=BookAmber,line width=0.8pt,rounded corners=1mm},
  external flow/.style={knowledge arrow,draw=BookMuted,line width=0.8pt,rounded corners=1mm},
  gradient flow/.style={knowledge arrow,draw=FigureLoss,line width=1.0pt,dashed,rounded corners=1mm},
  ema flow/.style={knowledge arrow,draw=BookMuted,line width=0.9pt,densely dotted,rounded corners=1mm},
  edge text/.style={inner sep=0.6mm,font=\figurefont\footnotesize,text=BookMuted,align=center}
]
  \draw[BookRule,line width=0.7pt] (1,-74) -- (167,-74);

  \node[anchor=north west,font=\figurefont\bfseries] at (1,1) {A\quad 获取与版本决策};
  \node[data node] (data) at (18,-14) {人工标签与\\待处理样本};
  \node[loop node] (model) at (56,-14) {当前模型\\版本$M^{(t)}$};
  \node[loop node] (candidate) at (94,-14) {候选预测\\及原始观测};
  \node[audit node] (selection) at (139,-14) {选择或复核\\接收训练目标};
  \node[state node] (input-change) at (18,-38) {预测请求修改\\调用安排};
  \node[loop node] (update) at (56,-38) {训练更新\\形成$M^{(t+1)}$};
  \node[audit node] (decision) at (94,-38) {继续、停止\\还是回退？};
  \node[audit node] (audit) at (139,-38) {开发验证\\过程审计};
  \node[state node] (previous) at (56,-61) {上一可用版本\\及旧结果};
  \node[state node,draw=FigureOutput] (freeze) at (94,-61) {冻结选定方案\\模型与预测请求};
  \node[state node,draw=FigureOutput,fill=FigureOutput!6] (holdout) at (139,-61) {留出验收\\独立质量证据};

  \draw[main flow] (data.east) -- (model.west);
  \draw[main flow] (model.east) -- (candidate.west);
  \draw[main flow] (candidate.east) -- (selection.west);
  \draw[evidence flow] (selection.south) -- (audit.north);
  \draw[evidence flow] (audit.west) -- (decision.east);
  \draw[evidence flow] (decision.west) -- node[edge text,above] {继续} (update.east);
  \draw[main flow] (update.north) -- (model.south);
  \draw[external flow] (input-change.east) -- (38,-38) -- (38,-14) -- (model.west);
  \draw[evidence flow] ([xshift=-6mm]decision.south) -- (88,-49) -| (previous.north);
  \node[edge text] at (72,-52) {回退};
  \draw[evidence flow] (decision.south) -- node[edge text,right] {停止} (freeze.north);
  \draw[main flow] (previous.east) -- (freeze.west);
  \draw[main flow] (freeze.east) -- (holdout.west);

  \node[anchor=north west,font=\figurefont\bfseries] at (1,-79) {B\quad Mean Teacher};
  \node[loop node,text width=22mm] (student) at (27,-95) {学生模型参数\\$\bm{\theta}^{(t)}$};
  \node[loop node,text width=22mm] (teacher) at (139,-95) {教师模型参数\\$\overline{\bm{\theta}}^{(t)}$};
  \node[state node,text width=22mm] (student-output) at (27,-114) {学生预测\\第$t$步};
  \node[loss node] (consistency) at (83,-114) {一致性损失\\$\mathcal{L}_{\mathrm{con}}^{(t)}$};
  \node[state node,text width=22mm] (teacher-output) at (139,-114) {教师目标\\不回传梯度};
  \node[loss node,text width=22mm] (supervised) at (27,-133) {监督损失\\$\mathcal{L}_{\mathrm{sup}}^{(t)}$};
  \node[data node,text width=12mm] (labels) at (61,-133) {真实标签};
  \node[loop node,text width=22mm] (optimizer) at (27,-153) {优化器\\学生梯度步};
  \node[loop node,text width=22mm] (student-next) at (70,-153) {学生模型参数\\$\bm{\theta}^{(t+1)}$};
  \node[loop node,text width=12mm] (ema) at (111,-153) {EMA};
  \node[loop node,text width=22mm] (teacher-next) at (149,-153) {教师模型参数\\$\overline{\bm{\theta}}^{(t+1)}$};
  \node[state node,text width=22mm,draw=FigureOutput,fill=FigureOutput!6] (next-target) at (149,-172) {重新预测\\第$t+1$步目标};

  \draw[main flow] (student.south) -- (student-output.north);
  \draw[main flow] (teacher.south) -- (teacher-output.north);
  \draw[main flow] (student-output.south) -- (supervised.north);
  \draw[main flow] (student-output.east) -- (consistency.west);
  \draw[main flow] (teacher-output.west) -- (consistency.east);
  \draw[main flow] (labels.west) -- (supervised.east);
  \draw[main flow] (student.west) -- (6,-95) |- (optimizer.west);
  \node[edge text,anchor=west] at (6.5,-124) {旧值};
  \draw[gradient flow] (supervised.south) -- (optimizer.north);
  \draw[gradient flow] (consistency.south) -- (83,-143) -| ([xshift=6mm]optimizer.north);
  \draw[main flow] (optimizer.east) -- (student-next.west);
  \draw[ema flow] (student-next.east) -- node[edge text,above] {$1-\alpha$} (ema.west);
  \draw[ema flow] (teacher.east) -- (164,-95) -- (164,-136) -| (ema.north);
  \node[edge text] at (137,-132) {旧教师模型参数，权重$\alpha$};
  \draw[ema flow] (ema.east) -- (teacher-next.west);
  \draw[main flow] (teacher-next.south) -- (next-target.north);
\end{tikzpicture}

  \caption{候选知识进入下一轮模型更新。面板A的开发验证与过程审计决定继续、停止或回退，选定方案冻结后另做留出验收。面板B实线是计算与参数读取，红色虚线是只更新学生的梯度，灰色点线是EMA。优化器使用旧学生参数与两项损失的梯度；EMA结合新学生与旧教师参数，再由新教师重新预测下一步目标。}
  \label{fig:knowledge-acquisition-iterative-update}
\end{figure*}

Mean Teacher中被匹配的是具有既定任务含义的预测，所以教师在本章提供任务目标。第~\ref{sec:data-derived-knowledge-cross-view-prediction}节的目标分支则产生表示向量，让另一分支去预测它，不先给出物种名称。二者都可能用EMA，差别在目标内容，不在更新公式或教师名称。硬伪标签和软一致性都可能继承初始错误，更新也不自动补足稀有类别；需同时监测接收规则、类别覆盖、人工参照标签上的任务结果及资源成本。

若采用整轮检查后替换教师，可参考Noisy Student：教师先生成未标注样本的伪标签，学生结合真实标签与伪标签训练，并在学生训练中加入噪声，再将学生升为下一轮教师\scite{knowledge-acquisition-noisy-student2020} 这里是版本替换，不是每步参数移动平均；新教师是否可用仍要经过开发检查，不能只凭完成了训练决定。

\subsection{输入内容的更新}
\label{sec:knowledge-acquisition-input-update}

若混淆来自不可见叶片，更新参数不能把叶片补到当前照片中。输入修改应先检查模型实际看过什么：类别边界是否写清，照片是否缺部位，示范是否少了反例，检索是否只匹配背景。这四项分别对应任务要求、样本上下文、示范和辅助证据，也限定哪些旧调用可能受修改影响。

确定成因以后，将输入变化保存为可比较的版本。作为构造例，规范$r_2$比$r_1$增加“叶缘锯齿和花期只能在对应部位或时间信息可见时使用”；示范$d_2$加入山茶与茶梅反例；样本$x_2$补入同株叶片；证据$e_2$换掉过时名录。每项都记录触发它的错误样本、适用对象与需要重算的旧调用。这样，“补充信息”就能落到具体内容与依赖，而不是覆盖原输入。

比较某项变化时，可固定模型、开发样本与评价口径，只替换一种输入成分，并对随机生成采用相同采样配置。若相对旧版本的物种错误减少且弃权仍符合规范，就有该输入变化在这些样本上有效的一项证据；仍须同时比较作答覆盖，避免只因作答变少而少犯错。补说明、换反例与加叶片分别检验规则表达、示范和观测的作用，但结果不能单独证明某项是唯一成因。反复针对同一小批错误修改，也会把输入设计拟合到这些样本；还须检查新批次、输入空间被挤占、说明冲突与额外检索成本。

模型与输入的变化应分开记录。两者同时更换，只能观察组合版本的结果，无法将差异分配给新参数或新证据；需要区分作用时可分阶段试验。完整方案连同模型、规范、上下文整理规则、示范和证据版本冻结后，才进入独立评价。第~\ref{chap:knowledge-quality-assessment-correction}章检验候选质量变化，本节将错误诊断落实为可追溯输入修改。

\subsection{获取安排的更新}
\label{sec:knowledge-acquisition-schedule-update}

版本修改以后，还须决定下一轮先处理谁。例如S042仍缺叶片，另一些照片已有可用答案但依赖旧模型，两类对象需要不同动作。\term{获取安排}（acquisition schedule）规定处理样本、调用次数、补证据或转人工的选择，以及全量还是局部重算。查询这些对象时，应分开内容、验证、执行阶段和版本依赖；“已确认可用”“需要重算”“证据不足”“等待人工”属于不同层，不能作互斥执行状态。S042可以同时物种未知、等待人工补观测。

执行状态沿用共同协议：尚未发起是未请求，已排队未开始是未处理，已开始未提交是处理中，操作完成是成功，异常终止未完成是失败。成功返回弃权仍是成功执行，失败则没有主动弃权动作。需要重算表示旧记录依赖已变化，不是第六种执行状态；为新请求另记进度与版本，并指定当前采用哪个结果，避免旧新答案都被当作现行记录。

重算范围应从变更及其依赖求出。整体参数更新，所有旧模型直接预测原则上都可能变化；仅改类别映射，可定位涉及该映射的记录；更新名录，要定位采用该来源或受其范围变化影响的候选；撤回近邻标签，要定位读取它的查询；改变图边，影响可能到达整片连通区域。只有依赖记录足以计算这些范围，才有局部刷新的依据，否则省下调用也会留下新旧版本不一致。

预算应与要补的缺口对应。新样本用于扩大覆盖，重复生成用于观察随机波动，困难复核用于查错，旧答案刷新用于处理版本变化，人工接管用于补入模型缺少的判断依据。相同模型、输入和证据若只反复产生相关答案，调用次数并不代表独立信息增加。表~\ref{tab:knowledge-acquisition-iteration-actions}据此列出反馈后的可选动作，而不是对每种现象规定唯一处理。

\begin{table*}[htbp]
  \centering
  \small
  \setlength{\RaggedRightParindent}{0pt}
  \renewcommand{\arraystretch}{1.15}
  \begin{tabularx}{\linewidth}{@{}>{\RaggedRight\arraybackslash}p{22mm}*{4}{>{\RaggedRight\arraybackslash}X}@{}}
    \toprule
    反馈情形 & 定位范围 & 可选动作 & 额外投入 & 过程检查 \\
    \midrule
    模型更新 & 旧参数或旧表示的依赖项 & 全量或局部刷新 & 训练、重建与推断 & 损失、覆盖与错误类型 \\
    规范变化 & 语义或字段变化的记录 & 重映射、重算或复核 & 规范评审与重算 & 边界案例与新增歧义 \\
    证据缺失 & 缺少必要信息的对象 & 补观测、检索或转人工 & 采集与来源核验 & 对象对应、弃权与错误 \\
    局部错误集中 & 特定类别、来源或图区域 & 定向复核与局部刷新 & 诊断、重建与复核 & 局部收益与其他区域退化 \\
    收益停滞 & 多轮重复的结果 & 停止、回退或换来源 & 比较与回退恢复 & 边际收益、覆盖与成本 \\
    \bottomrule
  \end{tabularx}
  \caption{反馈类型与可调整的获取安排。按反馈定位依赖范围，再比较可选动作、额外投入与过程结果；各行不是自动触发规则。开发验证和过程审计可支持继续、停止或回退，参与选择的数据不再作为冻结方案的最终独立验收依据。}
  \label{tab:knowledge-acquisition-iteration-actions}
\end{table*}

表中的定位还包括间接依赖。表示模型更新可能要求重建索引与图，规范变化可能连带改变类别映射和解析器，证据更新须回到原始来源及所用记录。局部修正以后，也要检查其他区域，因为参考关系或图边可能把变化传到原来没有错误的对象。

一批花卉照片因而可以分开安排：重算旧模型依赖的对象，给缺观测照片补拍叶片，把仍有歧义的对象转人工。若新版本只降低反复调试样本的错误，在新的开发批次中没有新增收益，就应结合成本选择停止或回退；输出变化变小也可能只是重复同一错误，不能证明质量提高。最终测试一旦用于这一选择，也就转成开发依据，须另留冻结后的独立验收。人工查询见第~\ref{chap:knowledge-acquisition-human}章，跨来源资源分配见第~\ref{chap:knowledge-fusion}章；本节确定的是模型流程内的执行顺序与版本一致性。

\section{本章小结}
\label{sec:knowledge-acquisition-model-based-summary}

从一批尚无物种答案的花卉照片出发，本章连起了模型准备、输入组织、候选计算与反馈修改。每条记录对应经过何种准备的模型、实际可见的任务要求与观测、所用示范或证据，以及产生答案的计算路径。复核再沿这些依赖决定改模型、补输入或重排获取。保存版本与范围，物种候选和弃权记录才有可追查、刷新或撤回的依据。

答案的依据随计算机制改变。直接预测读取参数中的任务关系，受训练覆盖与接口约束；近邻在调用时读取参考标签，依赖表示、距离与局部覆盖；图传播由任务标签锚定，借未知节点传递信息，也可能扩大错标签或错边的影响。三者可以组合，但不能因最终类别相同而省略中间依赖检查，也不能因数值收敛而认定候选正确。

候选回流训练时，任务教师提供目标，伪标签是暂时采用的任务答案，二者都没有新增正确性保证。Mean Teacher用参数移动平均更新教师，再由教师产生任务一致性目标；Self-Instruct则连同任务和输入一起生成材料，检查范围也从固定样本标签扩大到任务相容与材料分布。表示目标分支不先提供物种语义，不能仅因同用EMA就与任务教师等同。生成、筛选、参数更新和核验各回答不同问题。

是否继续迭代，要比较同一任务范围内的错误、弃权、覆盖与新增投入，而不是只看模型规模、输入长度、参考数或调用轮次。参与版本选择的数据提供开发依据，不能同时证明冻结方案的最终质量。本章得到的是可追溯的任务候选：多份结果怎样综合由第~\ref{chap:knowledge-fusion}章展开，可靠性与独立验收由第~\ref{chap:knowledge-quality-assessment-correction}章展开，从观测内部构造表示或分组目标则进入第~\ref{chap:data-derived-knowledge-acquisition}章。
